\documentclass[10pt,a4paper]{amsart}
\usepackage[margin=19mm]{geometry}
\usepackage{amsmath,amssymb,mathtools}
\usepackage[scr=rsfs,cal=euler]{mathalfa}
\usepackage{xfrac}
\usepackage[unicode,hidelinks,bookmarks=false]{hyperref}
\newcommand{\ie}{i.e.\ }
\newcommand{\cf}{cf.\ }
\newcommand{\resp}{respectively\ }
\newcommand{\Subsection}{\subsection}
\providecommand{\nobreakdash}{\nobreak\hbox{-}}
\setlength{\emergencystretch}{2em}
\multlinegap0pt
\usepackage{bm}
\usepackage{bbm}
\usepackage{dsfont}
\usepackage{xspace}
\usepackage{cancel}
\usepackage{mathtools}
\usepackage[shortlabels]{enumitem}
\usepackage{amsthm}
\makeatletter
\@ifundefined{lemma}{\newtheorem{lemma}{Lemma}}{}
\@ifundefined{proposition}{\newtheorem{proposition}[lemma]{Proposition}}{}
\@ifundefined{theorem}{\newtheorem{theorem}[lemma]{Theorem}}{}
\makeatother
\newcommand{\Aset}{\mathcal{A}}
\newcommand{\Bset}{\mathcal{B}}
\newcommand{\Xset}{\mathcal{X}}
\title[On the usage of $2$-node lines in $n$-correct and $GC\_n$ set]{On the usage of $2$-node lines in $n$-correct and $GC\_n$ sets}
\author{Hakop Hakopian, Gagik Vardanyan, Navasard Vardanyan}
\date{Source: arXiv:2508.13289v1 (CC BY 4.0)}
\begin{document}
\maketitle

\noindent\emph{Source and licence.} On the usage of $2$-node lines in $n$-correct and $GC_n$ sets. Version arXiv:2508.13289v1 (CC BY 4.0), distributed under the Creative Commons Attribution 4.0 International licence, as recorded in the arXiv licence metadata for this exact version and shown on the abstract page https://arxiv.org/abs/2508.13289v1. Full rights notice: licenses/arxiv-2508.13289-CC-BY-4.0.txt.\par\smallskip

\noindent\textit{Editorial scope.} Counted unit: the Theorem whose source label is \texttt{prp:m2ket} in the article (source0.tex lines 479--506, source label \texttt{prp:m2ket}), delivered together with the article's own complete original proof of it (source0.tex lines 507--581, one \texttt{proof} environment of 75 lines). No mathematics is written, repaired or re-derived: statement and proof are the article's own text line for line. Required context retained uncounted: maxmax (lines 276--276); maxcomp (lines 282--285); sptr2 (lines 341--342); th1sptr (lines 387--414). Every definition, display and label of those lines is retained unchanged. Omitted from this package and not redistributed: 1--275, 277--281, 286--340, 343--386, 415--478, 582--1093. Nothing inside a retained excerpt is omitted or altered: every retained body is delivered exactly as the article's lines are written, so no omission receipt of any kind accompanies this package. Citations: the retained excerpts carry no citation command at all, so no bibliography entry of the article is transcribed and no reference is redistributed. Reference adaptation: none was needed and none was made; every reference inside the delivered unit resolves inside it, so no unresolved reference or citation is produced. Printed slips kept literal: no printed word, symbol, hypothesis or number of the counted statement or of its proof was altered. Layout and class: the article's own class is \texttt{article} with packages including amssymb, amsthm, amsmath, graphicx, psfrag, array, hhline, enumitem, caption, subcaption; the wrapper is the lane's pinned \texttt{amsart} editorial wrapper at 10pt on A4, which restates the theorem-like environments the retained text needs and adds the article's own macro definitions that the retained text needs (\texttt{\textbackslash Aset}, \texttt{\textbackslash Bset}, \texttt{\textbackslash Xset}), with extra wrapper packages (bm, bbm, dsfont, xspace, cancel, mathtools, enumitem). The delivered numbering is the unit's own. Assets: no retained span contains a figure environment, a graphics-inclusion command, a PDF-page inclusion, a table or a TikZ picture; no image file is copied into this package, the package declares no asset dependency and no image file is redistributed with it. Licence: version arXiv:2508.13289v1 of this article is distributed under the Creative Commons Attribution 4.0 International licence, as recorded in the arXiv licence metadata for this exact version and shown on the abstract page https://arxiv.org/abs/2508.13289v1. A full rights notice is delivered at licenses/arxiv-2508.13289-CC-BY-4.0.txt. Characters: every retained excerpt is pure ASCII, so the delivered engine, XeLaTeX with its default Latin Modern text font, reports no missing character.
\begin{equation}\label{maxmax}q\in\Pi_k \hbox{ is a maximal curve}  \iff \Xset\setminus q  \hbox{ is an $(n-k)$-correct set.}\end{equation} 

\begin{proposition}\label{maxcomp} Let $\Xset$ be an $n$-correct set  and
a curve $q$ of degree $k\le n,$ be a maximal curve. Suppose that a curve $s$ of degree $m\le k,$ is a component of $q,$ that is $q=s r, \deg r=k-m.$
Then $s$ passes through  at least $d(n-k+m,m)$ nodes of $\Xset \setminus r.$ Moreover, $s$ passes through  exactly $d(n-k+m,m)$ nodes of $\Xset \setminus r$ if and only if $r$ is a maximal curve of degree $k-m.$ 
\end{proposition}

\begin{proposition} \label{sptr2} Let $\Xset$ be an $n$-correct set of nodes. Let also $\{A,B,C_1\}$ and $\{A,B,C_2\}$ be special triplets, where $A,B,C_1,C_2 \in\Xset$. Then $C_1=C_2.$
\end{proposition}

\begin{theorem}\label{th1sptr}
Let $\Xset$ be an $n$-correct set and $\mu:=\mu_k$ be a maximal curve of degree $k.$ Hence the set $\Bset:=\Bset_{n-k}=\Xset\setminus \mu$ is $(n-k)$-correct set. Denote by $\Aset:=\Aset_{k}=\Xset\cap \mu.$\\ Then a triplet of nodes $\{A,B,C\},$ where $A \in \Aset,\ B\in \Bset,$ and $C \in \Xset,$ is a special triplet if and only if the following two statements take place.
\begin{enumerate}
\setlength{\itemsep}{0mm}
\item
The line $\ell_1:=\ell_{AC}$ is a component of $\mu:\ \mu=\ell_1\mu_{k-1},$ in particular $B\notin \ell_1$ and $C\in \Aset.$  
\item
The set $\dot{\bar \ell}_1:=\{M\in{\ell_1}\cap  \Xset: \mu_{k-1}(M)\neq 0\}$ besides the nodes $A$ and $C,$ contains exactly $n-k$ more nodes of $\Xset,$ which 
coincide with the zeroes of $p^\star_{B,\Bset}\in\Pi_{n-k}$ on the line $\ell_1.$ Hence the latter zeroes are all real and distinct.
\end{enumerate}
Moreover, the statements \emph{(i) and (ii)} imply the following:
\begin{enumerate}
\setcounter{enumi}{2}
\setlength{\itemsep}{0mm}
\item
The curve $\mu_{k-1}:=\frac{1}{\ell_1}\ \mu$ is a maximal curve of degree $k-1.$ 

\item
The set $\Bset_{n-k+1}:=\Xset\setminus\mu_{k-1}=\Bset\cup \dot{\bar \ell}_1$ is an  $(n-k+1)$-correct set, for which $\ell_1$ is a maximal line.  
\item
$\Aset_{k-1}:=\Xset\cap\mu_{k-1}=\Aset\setminus \dot{\bar \ell}_1.$
\item 
$\{A,B,C\}$ is a special triplet in $\Bset_{n-k+1}$ too.
 
\item
If $\{B,D,E\}\subset\Xset$  is a special triplet, where $D\in \dot{\bar \ell}_1,$ then $\{B,D,E\}=\{A,B,C\},$ i.e., $\{D,E\}=\{A,C\}.$
\end{enumerate}
\end{theorem}

\begin{theorem}\label{prp:m2ket}
Let $\Xset, \mu:=\mu_k, \Bset:=\Bset_{n-k},\Aset:=\Aset_{k},$ be as in Theorem \ref{th1sptr}.
Let $\Xset, \mu:=\mu_k, \Bset:=\Bset_{n-k},\Aset:=\Aset_{k},$ be as in Theorem \ref{th1sptr}. Suppose that the distinct triplets $\{A_i,B,C_i\},\  1\le i\le m,$ with a common node $B$ are special, where $m\ge 2, A_i \in \Aset, B\in \Bset,$ and $C_i \in \Xset.$ 

Then, the following statements take place for each $i,\ 1\le i\le m:$
\begin{enumerate}
\setlength{\itemsep}{0mm}
\item
The line $\ell_i:=\ell_{A_iC_i}$ is a component of $\mu_{k-i+1},$ in particular  $B\notin\ell_i$ and $C_i\in\Aset_{k-i+1}.$  These $m$ components are distinct, hence $m\le k.$ Moreover, we have that
$$\{A_j,C_j\}\cap\ell_{A_iC_i} =\emptyset\ \ \forall\ 1\le j\le m,\ j\neq i.$$
\item
The set $\dot{\bar\ell}_i:=\{M\in{\ell_i}\cap  \Xset: \mu_{k-i}(M)\neq 0\}$ besides the nodes $A_i$ and $C_i,$ contains exactly $n-k+i-1$ nodes of $\Xset,$ from which $n-k$ coincide with the zeroes of $p^\star_{B,\Bset}\in\Pi_{n-k}$ on the line $\ell_i,$ which are all real and distinct, and $i-1$ are the following points of intersection: 
$$D_{ij}:=\ell_i\cap\ell_j\ \ \forall\ 1\le j \le i-1.$$ 
\item
The curve $\mu_{k-i}:=\frac{1}{{\ell_i}}\mu_{k-i+1}=\frac{1}{\ell_1\cdots\ell_i}\ \mu$  is a maximal curve of degree $k-i.$  
\item
The set  
$\Bset_{n-k+i}:=\Xset\setminus\mu_{k-i}=\Bset_{n-k+i-1}\cup \dot{\bar\ell}_i=\Bset\cup\dot{\bar\ell}_1\cup\cdots\cup\dot{\bar\ell}_i
$ 

is an  $(n-k+i)$-correct set, for which $\ell_i$ is a maximal line.
\item
$\Aset_{k-i}:=\Xset\cap\mu_{k-i}=\Aset_{k-i-1}\setminus\dot{\bar \ell}_i=\Aset\setminus (\dot{\bar \ell}_1\cup\cdots\cup\dot{\bar \ell}_i).$ 

\item
The line $\ell_i$ is a maximal line in  $\Bset_{n-k+s}$, where $i\le s\le m.$ In particular all the lines $\ell_1,\ldots,\ell_m$ are maximal lines in $\Bset_{n-k+m}.$
 \end{enumerate}
\end{theorem}

\begin{proof}
Let us use indiction on $m$ for the first two statements of Theorem.
The case $m=1$ follows from Theorem \ref{th1sptr}.
Let us assume that statements (i) and (ii) of Theorem are valid for $m$ and prove them for $m+1.$
Later we will prove, as in the case of Theorem \ref{th1sptr}, that statements (i) and (ii) imply the remaining statements (iii)-(vi). 

Let $\{A_i,B,C_i\},$ be the distinct special triplets $1\le i\le m+1,$ where $A_i \in \Aset,\ B\in \Bset,$ and  $C_i \in \Xset.$ 
Since the triplets are distinct, Proposition \ref{sptr2} readily implies that 
\begin{equation*}\label{ACij}\{A_i,C_i\}\cap \{A_j,C_j\}=\emptyset\ \ \forall\ 1\le i<j\le m+1.\end{equation*}

Now, let us check that $A_{m+1}\in \Aset_{k-m}.$

Indeed, in view of  Theorem \ref{th1sptr}, (vii), with $\{A,B,C\}=\{A_i,B,C_i\},\ 1\le i\le m,$ and $\{B,D,E\}=\{A_{m+1},B,C_{m+1}\},$ we conclude that $A_{m+1}\notin \dot{\bar\ell}_i,$ since $A_{m+1}$ differs from $A_i$ and $C_i.$  Hence (v), with $i=m,$ implies that $A_{m+1}\in\Aset_{k-m}.$  


Next, let us apply Theorem \ref{th1sptr} for the special triplet $\{A_{m+1},B,C_{m+1}\},$   the maximal curve $\mu_{k-m},$ and $\Aset_{k-m}:=\Xset\cap\mu_{k-m},\ \Bset_{n-k+m}:=\Xset\setminus\mu_{k-m}.$
Below we present the first two statements obtained and, after each one, we draw the corresponding conclusions.
\begin{enumerate}[label*=\arabic*.]
\setcounter{enumi}{0}
\setlength{\itemsep}{0mm}
\item
The line $\ell_{m+1}:=\ell_{A_{m+1}C_{m+1}}$ is a component of $\mu_{k-m},$ in particular $B\notin\ell_{m+1}$ and $C_{m+1}\in \Aset_{k-m}.$ 
\end{enumerate}
From here we conclude that the line $\ell_{m+1}$ differs from each of the lines
$\ell_1, \ldots,\ell_m.$ Indeed, the latter lines are components of $\mu:=\mu_k$ and if, say, $\ell_{m+1}=\ell_1,$ then we will have that $\ell_1$ is a component of $\mu_{k-m}=\frac{1}{\ell_1\cdots\ell_m}\ \mu.$ Hence $\ell_1$ is a multiple component of a maximal curve $\mu,$ which is contradiction.

Note that this proves (i) for $m+1,$ without the ``Moreover" part, since all the remaining statements of (i) hold by virtue of the induction hypothesis.
\begin{enumerate}[label*=\arabic*.]
\setcounter{enumi}{1}
\setlength{\itemsep}{0mm}
\item
The set $\dot{\bar\ell}_{m+1}:=\{M\in \ell_{m+1}\cap  \Xset: \mu_{k-m-1}(M)
\neq 0\}$ 
besides the nodes $A_{m+1}$ and $C_{m+1},$ contains exactly $n-k+m$ nodes of $\Bset_{n-k+m+1},$ from which $n-k$ are the zeroes of $p^\star_{B,\Bset}$ on the line $\ell_{m+1},$  and $m$ are the following points of intersection: 
$$D_{im+1}:=\ell_i\cap\ell_{m+1}\ \ \forall\ 1\le i \le m.$$ 
\end{enumerate}
Let us mention that the statement on $n-k+m$ nodes follows from Theorem \ref{th1sptr}, (ii), and  the formula
\begin{equation*}p^\star_{B,\Bset_{n-k+m}}=p^\star_{B,\Bset}\prod_{i=1}^m\ell_i,
\end{equation*}
 which is a direct consequence of Theorem \ref{th1sptr}, (iv).

\noindent This proves (ii). Note that the statements on the lines $\dot{\bar\ell}_i,\ 1\le i\le m,$ hold by virtue of the induction hypothesis.

Note also that above $n-k$ zeroes and $m$ points of intersection together make 
$n-k+m$ points. Since this number is the least possible (see the proof of Theorem \ref{th1sptr}), we conclude that all the mentioned
points are distinct. From here we get the ``Moreover" part of the statement (i).  

Next, let us verify that the statements (iii) to (vi) are valid.

First, in view of Proposition \ref{maxcomp}, with $m=1,$ we obtain that the curve 

$$\mu_{k-m-1}:=\frac{1}{\ell_{m+1}}\ \mu_{k-m}=\frac{1}{\ell_1\cdots\ell_{m+1}}\ \mu_{k}$$ is a maximal curve of degree $k-m-1.$ This gives (iii).
 

 Then, the relation \eqref{maxmax} implies that $\Bset_{n-k+m+1}$ is an  $(n-k+m+1)$-correct set, for which thus $\ell_{m+1}$ is a maximal line. The remaining equalities in (iv) and (v) follow from the definitions of $\mu_{k-m-1}$ and ${\dot{\bar{\ell}}}_{m+1}.$  Hence, the statements (iv) and (v) are valid. 

Finally, let us verify (vi). On the line $\ell_{i}$ we have $n-k+m+1$ nodes from $\Bset_{n-k+m}$ which can be divided into the following three groups:

$G_1:$ the nodes $A_i, C_i,$ 

$G_2:$ the $n-k$ nodes which are the zeroes of $p^\star_{B,\Bset}$ on the line $\ell_i,$  and 

$G_3:$ the $m-1$ intersection points: $\{D_{ij},\ \forall\ 1\le j\le m,\ j\neq i\}.$

Note that, according to (ii), these three sets are disjoint. 

For the first two groups,  in view of (ii) and (iv), we have that

$G_1\cup G_2\subset \dot{\bar\ell}_i\subset\Bset_{n-k+i}\subset\Bset_{n-k+m}.$

For the third group, similarly we obtain

$G_3\subset\dot{\bar\ell}_i\cup\cdots\cup\dot{\bar\ell}_s\subset\Bset_{n-k+m}.$
This proves (vi).
\end{proof}

\end{document}
